\chapter{Expected: errors as values}
\label{ch:expected}

\begin{aside}
Most C++ codebases handle errors with exceptions. \maya{} mostly
does not. The reason is the same reason \code{update} is pure: an
exception is a non-local control transfer, and non-local control
transfers undermine the value-typed contract the architecture
depends on. Where \maya{} needs to fail, it returns the failure
as a value. The tool is \code{maya::Expected<T, E>}: a sum type
that holds either a success or an error.
\end{aside}

\section{Why not exceptions?}

C++ exceptions are powerful. They unwind the stack, they call
destructors, they let you write code that ignores the error path
entirely. For application code with deep call graphs, they are
often the right tool.

For a library, they are often the wrong one. Specific issues:

\begin{itemize}[leftmargin=*]
  \item \textbf{Cost.} Throwing is expensive. Modern
    implementations (Itanium ABI on x86-64, MSVC's
    \emph{Funclets}) have nearly-zero cost on the happy path,
    but the cost on the failure path is large (table walks,
    destructor calls, allocator pressure). For a render loop
    that wants tens of thousands of frames per second on the
    happy path \emph{and} a fast crash-and-recover on the rare
    failure path, the asymmetry hurts.
  \item \textbf{Invisibility.} An exception can come out of
    almost any function. The reader does not know without
    reading the body --- or the bodies of every callee.
  \item \textbf{Composability with sum-typed code.} An exception
    bypasses the type system: it does not appear in the return
    type. Two pieces of code that handle their own failures
    through \code{Expected} compose cleanly; mix one piece that
    throws and the type system is no longer the source of truth.
  \item \textbf{Coroutines, restricted environments.} Some C++
    contexts (kernel modules, freestanding, \code{-fno-except})
    do not allow exceptions. Code that uses \code{Expected}
    works there.
\end{itemize}

For \maya{}'s low-level pieces --- terminal I/O, signal handling,
SIMD selection, file watching --- errors are values. Exceptions
are reserved for genuinely exceptional, unrecoverable conditions
(out of memory, hardware failure).

\begin{funfact}
Rust's \code{Result<T, E>} type is the same idea with a different
spelling. Haskell's \code{Either e a} is older still. Swift uses
\code{throws} but its model is closer to an exhaustively-checked
sum than to C++ exceptions. The pattern is older than C++.
\end{funfact}

\section{The Expected type}

\code{maya::Expected<T, E>} is a value-typed sum of a success
\code{T} and an error \code{E}:

\begin{lstlisting}
template <class T, class E>
class Expected {
    std::variant<T, E> v;
public:
    bool has_value() const { return v.index() == 0; }
    explicit operator bool() const { return has_value(); }
    T&       value()        { return std::get<0>(v); }
    const T& value() const  { return std::get<0>(v); }
    E&       error()        { return std::get<1>(v); }
    const E& error() const  { return std::get<1>(v); }
    // monadic helpers below
};
\end{lstlisting}

The C++23 standard introduced \code{std::expected<T, E>}; \maya{}
predates that, so it ships its own. The interface is similar.

\subsection*{Construction}

You can construct from either side:

\begin{lstlisting}
Expected<int, std::string> ok{42};                       // success
Expected<int, std::string> bad{std::unexpect, "boom"};   // error
\end{lstlisting}

The \code{std::unexpect} tag distinguishes the error constructor
from the success one when both types are convertible.

\subsection*{Checking}

\begin{lstlisting}
if (auto r = parse(input); r) {
    use(r.value());
} else {
    log(r.error());
}
\end{lstlisting}

The contextual bool conversion lets \code{if (r)} mean ``has
value''. Pair with the structured \code{auto r = \ldots; if (r)}
pattern to keep the scope tight.

\subsection*{Monadic helpers}

The interesting power is the chain operations:

\begin{itemize}[leftmargin=*]
  \item \code{transform(f)}: if value, apply \code{f}; if error,
    pass through. Returns a new \code{Expected<U, E>}.
  \item \code{and\_then(f)}: like \code{transform} but \code{f}
    itself returns an \code{Expected<U, E>}. Flattens the result.
  \item \code{or\_else(f)}: if error, apply \code{f} to the error
    (which may itself return an \code{Expected}); if value, pass
    through.
  \item \code{value\_or(default)}: extract the value or a default.
\end{itemize}

These let you write a pipeline that handles errors implicitly:

\begin{lstlisting}
Expected<Config, std::string> result =
    read_file("config.toml")
        .and_then(parse_toml)
        .and_then(validate)
        .transform(apply_defaults);

if (!result) std::cerr << result.error() << "\n";
\end{lstlisting}

If any stage fails, the rest are skipped and the error flows to
the end. No exceptions, no manual short-circuit ladders.

\begin{idea}
Monadic error handling is a style choice. Some prefer the
ladder; some prefer the chain. \code{Expected} supports both
and lets you pick per call site.
\end{idea}

\section{Where Expected appears in maya}

The main producers of \code{Expected} in \maya{}:

\begin{itemize}[leftmargin=*]
  \item \code{Terminal::open()} returns
    \code{Expected<Terminal, IoError>}. Failure cases:
    permission denied, device gone, raw-mode setup failed.
  \item \code{ansi::parse(byte\_stream)} returns
    \code{Expected<Event, ParseError>}. Failure cases:
    malformed escape, unterminated sequence.
  \item \code{file::read(path)} returns
    \code{Expected<std::string, IoError>}.
  \item \code{platform::detect\_simd()} returns
    \code{Expected<SimdLevel, std::string>}. Failure case: CPU
    feature detection failed.
  \item Widget initialisation in \file{maya/include/maya/widget/}
    --- e.g. \code{Picker::open(items)} returns
    \code{Expected<Picker, std::string>}.
\end{itemize}

The recurring shape: a function that interacts with the outside
world (OS, hardware, files) returns \code{Expected}. Pure
functions return their value directly.

\section{Error types}

\code{Expected} is parameterised on the error type. \maya{} uses
a few:

\begin{itemize}[leftmargin=*]
  \item \code{IoError}: a struct with a \code{std::errc} code, a
    \code{std::string} context message, and the syscall name.
  \item \code{ParseError}: a struct with a source offset and a
    description.
  \item \code{std::string}: when nothing more structured is
    needed.
\end{itemize}

You can use whatever error type fits. Tip: keep the error type
small and easy to construct, so the failure path is cheap.

\subsection*{IoError, in detail}

\begin{lstlisting}
struct IoError {
    std::errc code;          // standard error category
    std::string context;     // free-form, e.g. "opening /dev/tty"
    std::string syscall;     // "open", "read", "tcsetattr"
};
\end{lstlisting}

A constructor takes the \code{errno} value via
\code{std::make\_error\_condition} and stashes everything in one
place. When you log it:

\begin{lstlisting}
std::cerr << e.syscall << " failed ("
          << std::make_error_code(e.code).message() << "): "
          << e.context << "\n";
\end{lstlisting}

You get a message like ``open failed (No such file or
directory): /dev/tty12''.

\section{Crossing the boundary into Cmd}

In Chapter~\ref{ch:elm} we said \code{Cmd::task} returns a message.
What if the task fails?

The pattern is to make the message a sum:

\begin{lstlisting}
struct LoadOk    { std::string content; };
struct LoadFail  { IoError err; };
using Msg = std::variant<LoadOk, LoadFail, /* ... */>;

static Cmd<Msg> load_cmd(std::string path) {
    return Cmd<Msg>::task([path]() -> Msg {
        auto r = file::read(path);
        if (r) return LoadOk{std::move(r.value())};
        return LoadFail{std::move(r.error())};
    });
}
\end{lstlisting}

The task takes an \code{Expected}, peels it open, and routes the
result through the message variant. \code{update} then handles
both cases as branches. The error is data flowing through the
same pipeline as success.

\begin{warning}
Do not unwrap an \code{Expected} with \code{.value()} without
checking. \code{.value()} on an error throws. The whole point of
\code{Expected} is to avoid exceptions; calling \code{.value()}
without an \code{if (r)} undoes that.
\end{warning}

\section{Patterns}

\subsection*{Early return ladder}

\begin{lstlisting}
Expected<Config, std::string> load_config(std::string_view path) {
    auto raw = file::read(path);
    if (!raw) return std::unexpect, "read failed: " + raw.error().context;

    auto parsed = parse_toml(raw.value());
    if (!parsed) return std::unexpect, "parse failed: " + parsed.error();

    auto valid = validate(parsed.value());
    if (!valid) return std::unexpect, "validate failed: " + valid.error();

    return apply_defaults(valid.value());
}
\end{lstlisting}

The C-style ``check at every step'' approach. Verbose but
explicit; everyone can read it.

\subsection*{Monadic chain}

The same function with \code{and\_then}:

\begin{lstlisting}
Expected<Config, std::string> load_config(std::string_view path) {
    return file::read(path)
        .transform_error([](auto& e) {
            return "read failed: " + e.context;
        })
        .and_then(parse_toml)
        .and_then(validate)
        .transform(apply_defaults);
}
\end{lstlisting}

Shorter but requires the reader to understand the chain. Both
styles are fine; pick by audience.

\subsection*{Try-helper macro}

Some codebases define a \code{TRY} macro that hides the ladder:

\begin{lstlisting}
#define TRY(x) ({ auto __r = (x); if (!__r) return std::unexpect, __r.error(); std::move(__r.value()); })

Expected<Config, std::string> load_config(std::string_view path) {
    auto raw    = TRY(file::read(path));
    auto parsed = TRY(parse_toml(raw));
    auto valid  = TRY(validate(parsed));
    return apply_defaults(valid);
}
\end{lstlisting}

\maya{} does not use \code{TRY} in its public API (it relies on
a GCC extension), but the pattern is common in private code.

\begin{aside}
Rust has \code{?} for exactly this pattern, built into the
language. C++ is exploring similar syntax (P0779, ``Notes on the
syntax of try operators''), but it is years away. The macro is
the workaround.
\end{aside}

\section{Performance}

\code{Expected<T, E>} costs a variant: a tag byte plus enough
storage for the larger of \code{T} and \code{E}. For \code{T =
int} and \code{E = std::string}, that is around 40 bytes.

The happy path is one branch (the tag check); a missed branch
prediction is the worst case, perhaps a dozen cycles. The error
path is a string move, which is cheap.

Compared to exceptions:

\begin{itemize}[leftmargin=*]
  \item Happy path: \code{Expected} is one branch. Exceptions are
    zero branches (the unwinding tables sit dormant).
    Effectively a tie; the exception edges ahead by a hair on
    the very hot path.
  \item Error path: \code{Expected} is a value construction and a
    branch. Exceptions are a table walk, destructor calls,
    allocator pressure. \code{Expected} wins by an order of
    magnitude.
  \item Throughput: \code{Expected} is uniformly fast. Exceptions
    are bimodal.
\end{itemize}

For a renderer that wants predictable latency, \code{Expected} is
the better trade.

\section{Pitfalls}

\begin{warning}
\textbf{Bare types in the variant.} If \code{T} and \code{E}
are the same type, you cannot distinguish them. Use a wrapper
or a tag. \code{Expected<int, int>} compiles but is a foot-gun.
\end{warning}

\begin{warning}
\textbf{Implicit conversions from E to T.} If \code{T} accepts
construction from \code{E}, calling
\code{Expected<T, E>\{some\_E\}} might construct a success
holding a converted \code{T}. Use \code{std::unexpect} to be
unambiguous.
\end{warning}

\begin{warning}
\textbf{Discarded errors.} If you call a function returning
\code{Expected} and ignore the result, you lose the error. GCC
warns with \code{-Wunused-result} if the return type is marked
\code{[[nodiscard]]}; \maya{}'s \code{Expected} carries that
attribute.
\end{warning}

\begin{warning}
\textbf{Nested Expected.} \code{Expected<Expected<int, E>, E>} is
legal but ugly. Use \code{and\_then} to flatten, or refactor the
inner call to return the same error type.
\end{warning}

\section{When to throw}

Not every failure is recoverable. For genuinely exceptional
conditions --- assertion failures, out-of-memory, undefined
behaviour detected at runtime --- exceptions or \code{abort()}
are appropriate. \maya{} uses:

\begin{itemize}[leftmargin=*]
  \item \code{std::bad\_alloc}: propagated from allocations, not
    caught.
  \item \code{std::logic\_error}: thrown from APIs whose
    preconditions were violated (a sign of a bug in the caller).
  \item \code{abort()}: from internal sanity checks that should
    never fail (paired with an assert).
\end{itemize}

The dividing line: if the failure is something the caller might
reasonably want to recover from, use \code{Expected}. If the
failure indicates a bug, throw or abort.

\section*{Workshop}

\begin{enumerate}[leftmargin=*]
  \item Write a function \code{Expected<int, std::string>
    parse\_int(std::string\_view)} that returns the parsed
    integer or an error message.
  \item Chain three calls of \code{parse\_int} on different
    strings with \code{and\_then}, summing the results. What
    happens when one fails?
  \item Refactor a function in your code that throws an
    exception into one that returns \code{Expected}. Note where
    the type signature change improves the call site.
\end{enumerate}

\begin{recap}
\code{Expected<T, E>} is a value-typed sum of success and error.
It lets functions return failures as data, making the call sites
explicit and the type system the source of truth. \maya{} uses
it for every I/O boundary and reserves exceptions for genuinely
exceptional conditions.
\end{recap}
